\section{Related work and the taxonomy crosswalk}
\label{sec:related}

Seven prior works give the agent harness a structure, and they divide by what they measure.
Catalogues and completeness matrices list systems against a scheme, taxonomies derive a scheme from a
few systems read closely, and outcome studies compare harnesses by name. The closest to our design
question (RQ3) is \citet{guo2026qa2task}. Each group is limited for our purpose by a mechanism, not by
incompleteness.

\paragraph{Broader surveys of LLM agents}
Beyond these seven, surveys of LLM agents organize the field by agent component or method rather than
by harness. \citet{wang2023agentsurvey} propose a unified framework for constructing agents and review
applications and evaluation strategies, and \citet{xi2023riseagents} divide an agent into brain,
perception, and action. Later surveys decompose agents into components and orchestration patterns
\citep{agenticai2026taxonomies,agentsystems2026survey}, classify reasoning frameworks as single-agent,
tool-based, or multi-agent \citep{agenticreasoning2025survey}, treat workflows as computation graphs
fixed before or during a run \citep{workflowopt2026survey}, and frame post-training, memory, and skills
as adaptation of an agent or its tools \citep{adaptation2025survey}. Others cover one domain, software
engineering \citep{seagentic2025survey} or deep research \citep{deepresearch2025agents}, or one
artifact: \citet{agenticframeworks2025} compare seven frameworks, among them CrewAI, LangGraph, and
AutoGen, and \citet{skills2026lifecycle} reads a 124-paper audit set on skill libraries, assembled by
snowballing rather than a PRISMA-style query. The unit in these surveys is a paper, method, or
framework, not a sampled population of released systems, and none reports an inter-coder agreement
statistic. Outside the survey literature, \citet{frameworksfallshort2026} classify 5,669 bug reports
from five agent frameworks and find that 76.00\% manifest as incorrect functionality.

\paragraph{Catalogues and completeness matrices}
These works estimate what the field contains. \citet{li2026harnessengineering} apply the seven-layer
ETCLOVG taxonomy to a public catalogue of 171 entries at its 2026-05-08 snapshot (368 at our last
check), assigning one primary layer to each of 138 artifacts. The protocol borrows PRISMA's reporting
discipline~\citep{page2021prisma} but uses one primary coder and reports no kappa. Its unit also admits
sandboxes, tracers, gateways, and benchmarks, so its 368 and our corpus count different objects.
\citet{meng2026harnesssurvey} give a formal component definition, $H=(E,T,C,S,L,V)$, and rate 23
systems on it in a Harness Completeness Matrix. They call their method a structured expert survey rather
than a systematic review, and publish the matrix as a README table with no per-cell evidence. The
limitation is the cell itself. A label or rating records a verdict without the text it rests on, so a
reader cannot tell a component the system lacks from one its sources never mention. Without a second
coder, the rating also cannot be separated from the rater. Our coding keeps a quote behind every value
and gives silence a state of its own, which turns under-reporting into a measured rate.

\paragraph{Taxonomies from close reading}
These works give the strongest evidence per claim. \citet{rombaut2026insidescaffold} grew twelve
dimensions by open coding the source of 13 coding scaffolds at pinned commits. Every classification
carries a file path and line number, and a post-hoc pass checked 296 claims and corrected 19.
\citet{barbaste2026anatomy} read the source of eleven harnesses, one of them a non-coding contrast
point, at pinned releases, organize them into seven subsystems, and re-pin eight of them to observe
change over a quarter; they deliberately omit line-number citations. \citet{zhong2026responsibilities}
propose eleven component responsibilities and an H0--H3 ladder, a cumulative ablation that exposes more
of those responsibilities to one agent at each level while holding task, repository, and model fixed,
tested on a single constructed task. \citet{harnesscategorical2026} identifies the harness with a triple
$(G,\mathrm{Know},\Phi)$: wiring graph, structural certificates, and stage-to-model deployment map. Its
largest empirical test is a 10-instance SWE-bench-lite run in which no instance resolved. The limitation
is the cost of the evidence. Reading source code costs enough per system that both code studies stop at
eleven to thirteen systems. A sample that size cannot show how often a choice occurs, whether it is reported, or
whether it co-varies with outcomes, and Rombaut makes no claim about task success. We adopt Rombaut's
observation, classification, and evidence template and several of its value sets. We apply them to
documents rather than code, over 1,256 systems across domains, 247 of them coded twice.

\paragraph{Harness-identity outcome studies}
These works measure whether the harness matters. \citet{guo2026qa2task} decompose the harness into six
runtime responsibilities. Alone among the surveys, they compile same-model, different-harness evidence:
48 strictly matched Terminal-Bench 2.0 submissions and about 35 model--harness rows on
SWE-bench Verified. Among the 20 Terminal-Bench models with at least three harness results, the median
within-model range is 13.6 points, and 14 of 20 vary by at least 10. Comparisons that hold the model
fixed and swap the harness agree \citep{harnessbench2026,samemodel2026harness,scaffoldeffect2026}, and
\citet{zhang2026harnessreporting} argue that a comparison which does not disclose the harness cannot be
interpreted.
Under a fixed prompt, budget, and evaluator, changing only the harness moves Pass@1 by up to 27.4 points
\citep{clawswebench2026}, and a multi-agent scaffold raises SWE-bench Pro resolution by 6.3 to 18.5
points over three existing scaffolds with the same models \citep{icatagent2026}. Harnesses that edit
themselves improve held-out pass rates in all nine model--benchmark cells tested
\citep{selfharness2026}, the gains of evolved harnesses compensate recoverable execution defects
\citep{onerecipe2026harnesses}, and \citet{intentexecution2026} name the mismatch between what a model
intends and what its harness executes.
The limitation is the unit of contrast. A harness enters as a name, and a name bundles
dozens of design choices. A spread between two harnesses therefore has as many candidate causes as the
harnesses have differences: the design shows that the harness matters, not which choice does. Our
outcome analysis contrasts coded design dimensions within benchmark, split, and base-model keys, and
pools published within-study ablations by dimension. That step from identity to mechanism needs a coded
design, which Guo et al.\ set aside by choosing not to add another component taxonomy.

\begin{table}[t]
\centering
\scriptsize
\setlength{\tabcolsep}{3pt}
\caption{Crosswalk of our nine layers onto seven prior schemes. For the first five a mapping was
recorded before the schema froze (held in the replication package); Barbaste et al.\ and Zhong and
Zhu were mapped after the freeze from the full text. A parenthesis names the dimensions a slot
covers; the layer's other dimensions have no slot there. ``--'' marks no construct for the layer.
$^\dagger$A cross-cutting surface outside Barbaste et al.'s seven subsystems.}
\label{tab:crosswalk}
\begin{tabular}{@{}p{1.7cm}p{1.35cm}p{1.3cm}p{1.5cm}p{1.35cm}p{1.15cm}p{2.0cm}p{1.6cm}@{}}
\toprule
Our layer & Li et al.\ (ETCLOVG) & Guo et al. & Rombaut & Meng et al.\ $(E,T,C,$\allowbreak$S,L,V)$ & Banu &
Barbaste et al. & Zhong and Zhu \\
\midrule
A Context assembly & Context & observation, context & context strategy & $C$ (A2, A3) & -- &
LLM integration (A1); memory \& context (A2, A3) & context manager (A2) \\
B Tool interface & Tool interface & action & tool and environment interface & $T$ (B1--B5) & $G$ wiring &
tools \& actions (B1--B4); extensibility, orchestration (B5) & tool registry (B2) \\
C Control loop & Lifecycle and orchestration & control & control architecture & $E$ (C1) & $G$ wiring &
agent loop (C1, C2); orchestration (C3, C4); safety \& permissions (C5) & permission boundary (C5) \\
D Memory and state & Context & state & state management & $S$ (D1--D3) & coalgebraic state &
memory \& context (D1, D2); session substrate$^\dagger$ (D3) & task state (D1); project memory (D2) \\
E Verification and repair & Verification & verification & control architecture & $E$ (E2) & Know &
agent loop (E1, E2); memory \& context (E3) & verification protocol (E1) \\
F Budget and termination & Lifecycle and orchestration & control & resource management & $E$ (F1) & Know &
agent loop (F1--F3); LLM integration (F2) & -- \\
G Sandbox and environment & Execution environment & verification and governance (V) & execution
isolation & $L$ (G4) & -- & safety \& permissions (G1--G4) & permission boundary (G4) \\
H Observability and governance & Observability; Governance & not separated & -- & $V$ (H1, H3); $L$ (H4) &
Know & safety \& permissions (H4) & observability layer (H1); permission boundary (H4) \\
M Meta & -- & -- & -- & -- & $\Phi$ (M3) & LLM integration (M3) & -- \\
\bottomrule
\end{tabular}
\end{table}

Table~\ref{tab:crosswalk} shows that every prior scheme merges at least two of our layers. Li et al.\
place context assembly and memory in one layer, and control loop and budget in another; Guo et al.\
also merge control loop and budget, and Rombaut folds generate-test-repair verification into control
architecture. Down the columns, our observability and governance layer has no separate construct in Guo
et al.\ or Rombaut, and our sandbox layer has none in Banu. Under Meng et al.'s tuple, 14 of our 31
design dimensions (38 with metadata) have no slot, among them execution isolation, network policy, and replayability.
Li et al.\ and Banu each relate one pair of schemes, and \citet{barbaste2026anatomy} map Rombaut's
twelve dimensions onto their seven subsystems; Table~\ref{tab:crosswalk} sets all seven schemes side by side.

\citet{barbaste2026anatomy} merge three of our layers in their agent loop, which ``owns stop
conditions and failure recovery'' and so spans control loop, verification, and termination, and they
join context assembly with memory, and sandbox with guardrails, in their memory-and-context and
safety-and-permissions subsystems. None of their seven subsystems records tracing, replay, or
evaluation hooks, and three of their constructs have no dimension of ours: the extensibility
subsystem of hooks, skills, and plugins; the interface layer; and the provider-facing side of LLM
integration, such as streaming, reasoning effort, and per-model prompts, of which our schema keeps
only provider agnosticism (M3) and routing and caching as cost controls (F2).
\citet{zhong2026responsibilities} keep most of our layers apart but fold the permission model, action
guardrails, and approval gates (G4, H4, C5) into one permission boundary, and they have no construct
for our control loop beyond those gates and none for budget and termination. Four of their eleven
responsibilities have no dimension of ours: the task interface, failure attribution, the entropy
auditor, and the intervention logger.

The crosswalk also runs the other way, and there it exposes gaps in our schema. Three of Rombaut's
dimensions have no dimension of ours: loop driver, which that paper calls its most fundamental
distinction, tool discovery strategy, and multi-model routing. The last two are touched only partly.
\texttt{tool\_schema\_source} records whether a schema is hand-written or generated, not how
tools are registered, and \texttt{cost\_controls} records routing only as a cost mechanism.
Rombaut also separates memory the agent writes from read-only instruction files, which our
\texttt{long\_term\_memory} value \texttt{file\_notes} conflates. None of these gaps was repaired before
the schema froze; Section~\ref{sec:threats} weighs what they cost.

Four differences separate this review from the seven, each resting on an artefact a reader can inspect.
First, the review is pre-registered (\texttt{osf.io/ab2wn}) and every deviation is a numbered amendment.
Agreement is reported per dimension: on 247 double-coded systems, 37 of 38 dimensions reach a Cohen's
kappa of at least 0.6, and because both codings are model readings this is reproducibility, not correctness.
Second, the released matrix joins depth, breadth, and evidence: 1,256 systems on 38 dimensions, 47,728
cells, each valued cell carrying a quote and a locator. Third, the outcome analysis runs on coded design dimensions rather than
harness identity, over a comparable set published in full. Fourth, Table~\ref{tab:crosswalk} reconciles
seven prior schemes, five of them mapped before the schema froze and two after it.
